\section{Contexto y justificación del trabajo}
Las ciudades tienen una \textbf{deuda histórica con la diversidad humana} de sus habitantes. Hasta décadas recientes, el entorno urbano se ha \textbf{diseñado para la ciudadanía <<normativa>>}, definía como una persona adulta joven, ágil y sin ningún tipo de limitación física, sensorial o cognitiva. La irrupción del \textbf{vehículo privado} en el día a día de las personas añadió un nuevo elemento predominante en las vías urbanas (calzadas, aparcamientos). Esto eliminó espacios para los viandantes, facilitando entornos con \textbf{aceras estrechas, desniveles o problemas de continuidad} que excluyen a varios sectores de la sociedad en el ejercicio de sus derechos y disfrute de la ciudad.
\\

Cualquier persona puede pasar a lo largo de su vida de ser <<normativa>> a tener necesidades específicas de movilidad. Por esa razón, la \textbf{\gls{accesibilidad-universal}} no es una demanda exclusiva de la ciudadanía con discapacidad permanente, sino una condición que \textbf{afecta a toda la sociedad} en distintas etapas de la vida.
\\

\begin{figure}[H]
  \centering
  \begin{tikzpicture}[font=\sffamily, x=1cm, y=1cm]
    % --- Bloque de identidad con el titular
    \fill[uoc-cian] (0,0) rectangle (3.8,5.2);
    \node[anchor=north west, text width=3.3cm, align=left, text=uoc-azul,
          font=\sffamily\large\bfseries] at (0.3,4.8)
      {\begin{hyphenrules}{nohyphenation}¿A quién beneficia la accesibilidad?\end{hyphenrules}};
    \node[anchor=south west, text width=3.1cm, align=left, text=tinta,
          font=\sffamily\scriptsize] at (0.3,0.3) {España};
 
    % --- Filetes del panel (como la cabecera UOC)
    \draw[uoc-cian-filete, line width=1mm] (3.8,5.15) -- (16.7,5.15);
    \draw[uoc-cian-filete, line width=1mm] (3.8,0.05) -- (16.7,0.05);
 
    % --- Tres cifras: \x = inicio de columna, \pct = % para la minibarra
    \foreach \x/\cifra/\pct/\txt/\base in {
        4.3/{9,5\,\%}/9.5/{de la población tiene una \textbf{discapacidad reconocida}: 4,38 millones de personas}/{Base: población total},
        8.4/{\textgreater\,30\,\%}/30/{de la población es \textbf{beneficiaria real} de la accesibilidad universal}/{Base: población total},
        12.5/{32,3\,\%}/32.3/{de las personas \textbf{mayores de 65 años} tiene alguna discapacidad}/{Base: mayores de 65 años}} {
      % cifra
      \node[anchor=north west, inner sep=0pt, text=uoc-azul, font=\sffamily\bfseries, scale=2.2]
        at (\x,4.75) {\cifra};
      % texto
      \node[anchor=north west, inner sep=0pt, text width=3.6cm, align=left, text=tinta,
            font=\sffamily\small] at (\x,3.5) {\begin{hyphenrules}{nohyphenation}\txt\end{hyphenrules}};
      % minibarra sobre el 100 % de su base
      \fill[rejilla] (\x,1.05) rectangle ++(3.6,0.3);
      \fill[uoc-azul] (\x,1.05) rectangle ++(3.6*\pct/100,0.3);
      \node[anchor=north west, inner sep=0pt, text=tinta-suave, font=\sffamily\scriptsize]
        at (\x,0.85) {\base};
    }
    % separadores verticales finos
    \foreach \x in {8.1,12.2} \draw[filete, line width=\trazofino] (\x,0.5) -- (\x,4.7);
  \end{tikzpicture}
  \caption{La accesibilidad en tres cifras.}
  \label{fig:infografia_accesibilidad}
  \fuente{INE, Encuesta de Discapacidad, Autonomía personal y situaciones de Dependencia (EDAD) 2020; Observatorio 2030 CSCAE y Fundación ONCE, \emph{Documenta 1.2}.}
\end{figure}




En España, según información de 2020, \textbf{4,38 millones de habitantes tenían algún tipo de discapacidad reconocida} (física, visual, auditiva, cognitiva o relacional), lo que representa aproximadamente un 9,5\% de la población. Este porcentaje se va acentuando con la edad, alcanzando valores superiores al \textbf{30\% en las personas mayores de 65 años} (colectivo que representa un 20\% del total y en constante aumento por el envejecimiento demográfico). No se pueden olvidar tampoco las personas con \textbf{\gls{movilidad-reducida-temporal}} (lesiones, uso de muletas o silla de ruedas) ni las \textbf{familias con carritos de bebé}.
\\

Todos estos colectivos representan más de un 30\% de la ciudadanía, y reflejan que la problemática de la accesibilidad no solo impacta en ciertos colectivos, sino en amplios grupos de personas en el momento actual, y que \textbf{nos afectará en mayor o menor medida a todos}.
\citadestacada{La accesibilidad es imprescindible para el 10\,\% de la población, necesaria para el 40\,\% y cómoda para el 100.}{PREDIF}

De esta necesidad, y con el objetivo de tratar de tirar abajo la deuda accesible de las ciudades, se han redactado varias leyes en la legislación estatal española, de las comunidades autónomas y municipal.
El marco estatal de referencia es la \textbf{\gls{orden-tma-851-2021} de 23 de julio}, que legisla sobre las \textbf{condiciones básicas de accesibilidad} y no discriminación, y los derechos de utilización de los espacios públicos urbanos.
\\

\begin{figure}[htbp]
  \centering
  \begin{tikzpicture}[tfm, font=\sffamily\scriptsize,
      every node/.append style={execute at begin node={\begin{hyphenrules}{nohyphenation}}, execute at end node={\end{hyphenrules}}},
      cota/.style={{Stealth[length=3.5pt]}-{Stealth[length=3.5pt]}, draw=uoc-azul, line width=0.6pt},
      cotatxt/.style={font=\sansmath\sffamily\scriptsize\bfseries, text=uoc-azul, fill=papel, inner sep=1.2pt},
      ref/.style={draw=uoc-azul, line width=\trazofino, densely dotted},
      mobiliario/.style={fill=tinta-suave, draw=tinta},
      rotulo/.style={font=\sffamily\scriptsize, text=tinta-suave},
    ]
    % =============== a) PLANTA (vista cenital, como en la ortofoto) ===============
    \begin{scope}[x=1.15cm, y=1.15cm]
      \node[anchor=south west, font=\sffamily\small\bfseries, text=tinta] at (0,4.45) {Itinerario peatonal accesible: dimensiones mínimas};
      % edificación
      \fill[rejilla] (0,3.2) rectangle (12,4.2);
      \draw[tinta-suave, line width=\trazo] (0,3.2) -- (12,3.2);
      \node[rotulo] at (6,3.7) {Edificación · línea de fachada};
      % acera
      \fill[papel] (0,0) rectangle (12,3.2);
      % IPA: banda libre junto a la fachada
      \fill[uoc-cian-suave] (0,1.4) rectangle (12,3.2);
      \draw[uoc-azul, dashed, line width=\trazo] (0,1.4) -- (12,1.4);
      \node[font=\sffamily\small\bfseries, text=uoc-azul] at (5.2,2.5) {Itinerario peatonal accesible (IPA)};
      \node[font=\sffamily\scriptsize, text=uoc-azul] at (5.2,2.05) {continuo y sin obstáculos};
      % banda exterior (mobiliario)
      \node[rotulo, anchor=west] at (9.35,0.7) {Banda exterior: mobiliario};
      % bordillo y calzada
      \fill[seg-calzada] (0,-1.8) rectangle (12,0);
      \draw[tinta, line width=1.6pt] (0,0) -- (12,0);
      \draw[papel, line width=1pt, dash pattern=on 10pt off 8pt] (0,-1.1) -- (12,-1.1);
      \node[font=\sffamily\scriptsize, text=papel] at (1.3,-0.55) {Calzada};
            % mobiliario a >= 0,40 m del bordillo
      \fill[mobiliario] (2.4,0.7) circle (0.15);            \node[rotulo, below=1pt] at (2.4,0.55) {farola};
      \filldraw[mobiliario] (3.7,0.5) rectangle (4.1,0.9);   \node[rotulo, below=1pt] at (3.9,0.5) {papelera};
      \draw[tinta, fill=rejilla] (5.9,0.4) rectangle (6.9,1.2);
      \fill[seg-vegetacion] (6.4,0.8) circle (0.28);         \node[rotulo, below=1pt] at (6.4,0.4) {alcorque};
      \foreach \x in {8.2,8.8} \fill[mobiliario] (\x,0.55) circle (0.09);
      \node[rotulo, below=1pt] at (8.5,0.45) {bolardos};
      % cotas
      \draw[ref] (11.2,1.4) -- (11.6,1.4);
      \draw[cota] (11.4,1.4) -- (11.4,3.2) node[midway, cotatxt, rotate=90] {$\geq$\,1,80\,m};
      \draw[ref] (2.4,0.55) -- (1.55,0.55);
      \draw[cota] (1.65,0) -- (1.65,0.55) node[midway, cotatxt, left=1pt] {$\geq$\,0,40\,m};
      \node[rotulo, anchor=north east] at (12,-0.02) {\textcolor{papel}{bordillo}};
      % escala gráfica
      \begin{scope}[shift={(9.6,-1.55)}]
        \fill[papel] (0,0) rectangle (1,0.08); \fill[tinta-suave] (1,0) rectangle (2,0.08);
        \draw[papel, line width=0.4pt] (0,0) rectangle (2,0.08);
        \node[font=\sffamily\tiny, text=papel, above] at (0,0.08) {0};
        \node[font=\sffamily\tiny, text=papel, above] at (2,0.08) {2\,m};
      \end{scope}
    \end{scope}

  \end{tikzpicture}
  \caption{Itinerario peatonal accesible: dimensiones mínimas en plano.}
  \label{fig:ipa-plano}
  \fuente{elaboración propia a partir de la Orden TMA/851/2021.}
\end{figure}


En esta normativa se impone como mínimo a las administraciones municipales:
\begin{itemize}
    \item La obligación de elaborar los \textbf{\gls{PAU}}, planes estratégicos para identificar el estado de las calles, programar y priorizar las intervenciones y comprobar su cumplimiento.
    \item Impulsar el concepto de \textbf{accesibilidad universal}, no solo pensando en la población con problemas de movilidad, sino atendiendo a todos los tipos de diversidad.
    \item Establecer los \textbf{requisitos mínimos} para considerar las vías urbanas como accesibles, bajo el concepto de \textbf{Itinerario Peatonal Accesible (IPA)}.
\end{itemize}
 
\begin{figure}[htbp]
  \centering
  \begin{tikzpicture}[tfm, font=\sffamily\scriptsize,
      every node/.append style={execute at begin node={\begin{hyphenrules}{nohyphenation}}, execute at end node={\end{hyphenrules}}},
      cota/.style={{Stealth[length=3.5pt]}-{Stealth[length=3.5pt]}, draw=uoc-azul, line width=0.6pt},
      cotatxt/.style={font=\sansmath\sffamily\scriptsize\bfseries, text=uoc-azul, fill=papel, inner sep=1.2pt},
      ref/.style={draw=uoc-azul, line width=\trazofino, densely dotted},
      mobiliario/.style={fill=tinta-suave, draw=tinta},
      rotulo/.style={font=\sffamily\scriptsize, text=tinta-suave},
    ]
    % =============== b) SECCIÓN TRANSVERSAL ===============
    \begin{scope}[shift={(0,-7.4)}, x=1.15cm, y=1.15cm]
      \node[anchor=south west, font=\sffamily\small\bfseries, text=tinta] at (0,3.35) {Vista transversal del itinerario peatonal accesible};
      % calzada (más baja) y acera
      \fill[seg-calzada] (0,-0.35) rectangle (1.5,-0.15);
      \fill[rejilla] (1.5,-0.35) rectangle (4.7,0);
      \draw[tinta, line width=\trazo] (0,-0.15) -- (1.5,-0.15) -- (1.5,0) -- (4.7,0);
      % fachada
      \fill[rejilla] (4.7,-0.35) rectangle (5.1,3.1);
      \draw[tinta-suave, line width=\trazo] (4.7,0) -- (4.7,3.1);
      % volumen libre del IPA
      \fill[uoc-cian-suave] (2.9,0) rectangle (4.7,2.2);
      \draw[uoc-azul, dashed, line width=\trazo] (2.9,0) -- (2.9,2.2) -- (4.7,2.2);
      % persona genérica
      \fill[uoc-azul] (3.8,1.62) circle (0.13);
      \fill[uoc-azul, rounded corners=2pt] (3.62,0.02) rectangle (3.98,1.45);
      \node[font=\sffamily\scriptsize\bfseries, text=uoc-azul] at (3.3,1.95) {IPA};
      % farola en la banda exterior
      \fill[tinta-suave] (1.95,0) rectangle (2.03,2.9);
      \fill[tinta-suave] (2.03,2.84) rectangle (2.6,2.9);
      % elemento volado (toldo o banderola) por encima de 2,20 m
      \fill[tinta-suave] (4.7,2.45) -- (3.9,2.35) -- (3.9,2.3) -- (4.7,2.3) -- cycle;
      % cotas
      \draw[cota] (2.9,-0.55) -- (4.7,-0.55) node[midway, cotatxt, below=1pt, fill=none] {$\geq$\,1,80\,m};
      \draw[ref] (2.9,0) -- (2.9,-0.65); \draw[ref] (4.7,-0.35) -- (4.7,-0.65);
      \draw[cota] (2.65,0) -- (2.65,2.2) node[midway, cotatxt, rotate=90] {$\geq$\,2,20\,m};
      \draw[ref] (2.55,2.2) -- (2.9,2.2);
      \draw[cota] (1.5,0.35) -- (1.95,0.35) node[midway, cotatxt, above=1pt, fill=none] {$\geq$\,0,40\,m};
      \node[rotulo, align=center] at (0.75,0.3) {calzada};
    \end{scope}
 
    % =============== Requisitos ===============
    \begin{scope}[shift={(7.0,-7.4)}]
      \node[anchor=north west, text width=7.2cm, align=left, font=\sffamily\scriptsize, text=tinta,
            draw=uoc-azul, line width=\trazofino, fill=superficie, inner sep=6pt] at (0,3.75) {%
        {\small\bfseries\color{uoc-azul} Requisitos del IPA}\\[2pt]
        {\color{tinta-suave}Orden TMA/851/2021}\\[5pt]
        \textbf{Ancho libre de paso} $\geq$\,1,80\,m, continuo y sin obstáculos.\\[3pt]
        \textbf{Altura libre de paso} $\geq$\,2,20\,m en todo el ancho del itinerario.\\[3pt]
        \textbf{Mobiliario urbano} en la banda exterior de la acera, a $\geq$\,0,40\,m del borde del bordillo.\\[3pt]
        \textbf{Trazado} junto a la línea de fachada o elemento que la sustituya.};
    \end{scope}
  \end{tikzpicture}
  \caption{Itinerario peatonal accesible: dimensiones mínimas en sección transversal.}
  \label{fig:ipa-seccion}
  \fuente{elaboración propia a partir de la Orden TMA/851/2021.}
\end{figure}



A pesar de esta legislación, según los diagnósticos realizados por el IMSERSO, el estado actual de la accesibilidad de las ciudades es bastante negligente. En el año 2005, más del \textbf{16\% de aceras eran totalmente inaccesibles} por tener anchos inferiores a lo establecido por ley, y peor aún, más del \textbf{63\% de los itinerarios peatonales tenían obstáculos físicos} que dificultan la movilidad de los viandantes. Según el informe de Índice de Accesibilidad e Inclusión de la Fundación ONCE (2022), el \textbf{37\%} de las necesidades de accesibilidad en la vía urbana \textbf{siguen incumpliendo la ley}.
\\

La ejecución de los planes de accesibilidad sufre de unos procedimientos excesivamente lentos y costosos, debido a que las inspecciones del equipo técnico deben hacerse de manera presencial y con materiales específicos, lo que conlleva un:
\begin{itemize}
    \item \textbf{Coste económico elevado:} por la mano de obra, los desplazamientos y el uso de material, especialmente en localidades más pequeñas.
    \item \textbf{Rápida caducidad de los diagnósticos:} al ser tan lentos, no se adaptan al crecimiento dinámico de las ciudades, como obras o cambios en el mobiliario urbano.
    \item \textbf{Procesamiento lento de los datos:} una vez obtenidas las mediciones, los datos se tienen que subir a las plataformas web, en algunos casos inexistentes.
\end{itemize}


La necesidad de acelerar los planes de accesibilidad es el punto de partida de este proyecto. Para ello, se plantea la utilización de parámetros potencialmente obtenibles a partir de imágenes aéreas y a pie de calle, así como datos cartográficos públicos, como fuente de datos para el diagnóstico de las vías urbanas, utilizando técnicas de inteligencia artificial, construyendo un \textbf{modelo de clasificación de accesibilidad} la evaluación de las calles de cada ciudad.
\\

Este planteamiento permitirá \textbf{reducir la carga de trabajo, la lentitud y los costes} que conlleva la inspección \textit{in situ} del procedimiento actual, \textbf{impactando} positivamente en el conjunto de la \textbf{ciudadanía}.
\\

\begin{figure}[htbp]
  \centering
  \begin{tikzpicture}[tfm,
      every node/.append style={execute at begin node={\begin{hyphenrules}{nohyphenation}}, execute at end node={\end{hyphenrules}}}, % sin guiones
      paso/.style={arqbase, text width=22mm, minimum height=11mm, inner sep=3pt,
                   font=\sffamily\scriptsize, align=center},
      manual/.style={paso, fill=superficie, draw=tinta-suave},
      datos/.style={paso, fill=arq-datos},
      algoritmo/.style={paso, fill=nn-conv, double=nn-conv, double distance=1.2pt},
      externo/.style={paso, fill=arq-externo, draw=acc-deficiente},
      % Etapas del objetivo secundario (línea discontinua)
      extension/.style={paso, fill=nn-conv!40, draw=uoc-azul, dashed},
      extext/.style={paso, fill=arq-externo, draw=acc-deficiente, dashed},
      nota/.style={font=\sffamily\scriptsize, text=tinta, anchor=north west, inner sep=0pt,
                   text width=41mm, align=left},
    ]
    % ---------- Carriles (fondo con mayor altura)
    \fill[uoc-gris]       (-2.25, 2.00) rectangle (12.45,-2.00);
    \fill[uoc-cian-suave] (-2.25,-2.30) rectangle (12.45,-7.10);

    % Etiquetas laterales a dos líneas (la segunda línea no negrita) para que quepan en el carril
    \node[rotate=90, align=center, font=\sffamily\scriptsize, text=tinta-suave] at (-1.8, 0.0)  {\textbf{ACTUAL}\\manual (\emph{in situ})};
    \node[rotate=90, align=center, font=\sffamily\scriptsize, text=uoc-azul]    at (-1.8,-4.65) {\textbf{PROPUESTO}\\variables seleccionadas};

    % ---------- Proceso actual
    \node[manual] (a1) at (0,0)    {Planificación de rutas por barrios o sectores};
    \node[manual] (a2) at (2.75,0)  {Desplazamiento del equipo técnico};
    \node[manual] (a3) at (5.5,0)  {Medición manual calle a calle};
    \node[manual] (a4) at (8.25,0)  {Registro en informes y fotografías};
    \node[datos]  (a5) at (11,0) {Carga manual en la plataforma};
    \foreach \i/\j in {a1/a2,a2/a3,a3/a4,a4/a5} \draw[flecha] (\i) -- (\j);
    \draw[flecha, dashed] (a3.north) .. controls +(0,8mm) and +(0,8mm) .. (a2.north)
      node[pos=0.5, above, etiqueta] {se repite en cada calle};
    \foreach \x/\t in {-1.1/{Coste elevado de personal y desplazamientos},
                        3.45/{Foto puntual del diagnóstico, se queda obsoleta pronto},
                        8.0/{Datos lentos de procesar, difícil tener una visión global}}
      \node[nota] at (\x,-1.0) {\textcolor{estado-bloqueado}{\textbf{\textminus}}~\t};

    % ---------- Proceso propuesto
    % Etapas discontinuas = objetivo secundario (extracción de variables desde imagen).
    % En el trabajo principal las variables se toman de los datos de campo de los ayuntamientos.
    \node[extext]    (b1) at (0,-4.3)    {Imágenes aéreas y a pie de calle};
    \node[extension] (b2) at (2.75,-4.3)  {Extracción de variables seleccionadas};
    \node[datos]     (b3) at (5.5,-4.3)  {Variables seleccionadas por tramo};
    \node[algoritmo] (b4) at (8.25,-4.3)  {Modelo de clasificación de accesibilidad};
    \node[datos]     (b5) at (11,-4.3) {Resultados y datos exportables};
    \foreach \i/\j in {b1/b2,b2/b3,b3/b4,b4/b5} \draw[flecha] (\i) -- (\j);
    \node[externo, text width=34mm, minimum height=8mm] (v) at (8.25,-2.95) {Datos de campo de Albacete y Algeciras};
    \draw[flecha, dashed] (v) -- (b4);
    \draw[flecha, dashed] (b5.south) .. controls +(0,-8mm) and +(0,-8mm) .. (b1.south)
      node[pos=0.5, below=1pt, etiqueta, fill=uoc-cian-suave] {se actualiza con cada nueva carga de imágenes};

    % Términos positivos desplazados a -6.20
    \foreach \x/\t in {-1.1/{Cobertura más amplia con menos visitas de campo},
                        3.45/{Se identifica qué variables bastan para clasificar},
                        8.0/{Facilita la priorización de intervenciones}}
      \node[nota] at (\x,-6.20) {\textcolor{estado-completado}{\textbf{+}}~\t};
  \end{tikzpicture}

  \vspace{3mm}
  % ---------- Leyenda
  % Los elementos se encadenan con posiciones relativas (no coordenadas fijas),
  % de modo que no se solapan aunque la tipografía sea más ancha.
  \begin{tikzpicture}[font=\sffamily\scriptsize, text=tinta,
      leg/.style={minimum width=4.5mm, minimum height=2.4mm, inner sep=0pt, anchor=west, line width=\trazo},
      legt/.style={anchor=west, inner sep=0pt}]
    \node[leg, draw=tinta-suave, fill=superficie] (k1) at (0,0) {};
    \node[legt] (t1) at ([xshift=1.2mm]k1.east) {Tarea manual};
    \node[leg, draw=uoc-azul, fill=nn-conv, double=nn-conv, double distance=0.8pt] (k2) at ([xshift=3.5mm]t1.east) {};
    \node[legt] (t2) at ([xshift=1.2mm]k2.east) {Algoritmo};
    \node[leg, draw=uoc-azul, fill=arq-datos] (k3) at ([xshift=3.5mm]t2.east) {};
    \node[legt] (t3) at ([xshift=1.2mm]k3.east) {Datos / resultado};
    \node[leg, draw=acc-deficiente, fill=arq-externo] (k4) at ([xshift=3.5mm]t3.east) {};
    \node[legt] (t4) at ([xshift=1.2mm]k4.east) {Fuente externa};
    \node[leg, draw=uoc-azul, fill=nn-conv!40, dashed] (k5) at ([xshift=3.5mm]t4.east) {};
    \node[legt] (t5) at ([xshift=1.2mm]k5.east) {Extensión};
    \draw[flecha, dashed] ([xshift=3.5mm]t5.east) -- ++(5mm,0) coordinate (k6);
    \node[legt] (t6) at ([xshift=1.2mm]k6) {Repetición};
  \end{tikzpicture}
  \caption{Proceso actual de diagnóstico de la accesibilidad mediante inspección \emph{in situ} frente al proceso propuesto basado en variables seleccionadas.}
  \label{fig:proceso_actual_propuesto}
  \fuente{elaboración propia, con ayuda de Claude y Gemini.}
\end{figure}


\section{Motivación personal}
Mi interés personal a la hora de seleccionar este trabajo parte principalmente de dos motivos. En primer lugar, buscaba realizar un proyecto con\textbf{ carga social} y que tuviera \textbf{aplicación real y humana}, intentando impactar en la mejora de la vida de las personas.
\\

En segundo lugar, quería que estuviese \textbf{relacionado con la accesibilidad}, debido a que en mi \textbf{círculo más cercano} cuento con \textbf{personas con diversidad funcional}, como la hipoacusia o problemas de movilidad, y gracias a su activismo y su relación con plataformas de accesibilidad como \textit{PREDIF} o \textit{Music For All}, me ha permitido conocer más de cerca el trabajo que realizan estas \textbf{asociaciones}, y el largo camino que queda por recorrer, además de conocer gente con diferentes diversidades.
\\

Una vez solicitado el proyecto, tuve la mala suerte de \textbf{sufrir una lesión física} en el pie, que me ha obligado durante semanas a \textbf{utilizar silla de ruedas y muletas para desplazarme}.
\\

Esta vivencia personal me ha permitido \textbf{comprender mucho mejor todos los obstáculos} con los que las ciudades nos impiden desplazarnos con autonomía e independencia.
Además, he podido comprobar que la \textbf{accesibilidad} no es algo estático, sino \textbf{dinámico y dependiente del tiempo}, lo que hace que calles en principio accesibles se transformen en problemáticas.
\\

Un ejemplo es la presencia de cubos de basura depositados en la acera por la noche en el centro de Madrid, que reducen de manera drástica el ancho de las aceras, la disposición de mobiliario urbano, como las marquesinas de autobús de manera que se impida el paso incluso para los viandantes; o el \textbf{incivismo de la sociedad}.
\\

Esta experiencia personal me permitirá ayudar a identificar características y variables a la hora de ejecutar el proyecto.
\\

\section{Objetivos del trabajo}

\subsection{Objetivo general}

El objetivo general del trabajo es el \textbf{diseño e implementación de un modelo de clasificación de parámetros de accesibilidad del viario urbano} basado en técnicas de \gls{machine-learning}, a partir de conjuntos de datos de estudios de accesibilidad de dos ciudades (Albacete y Algeciras).
Se tratará de evaluar la viabilidad de clasificar la accesibilidad de las calles, a partir de un subconjunto de aquellos parámetros que se puedan obtener a partir de imágenes aéreas de alta resolución, o imágenes obtenidas desde otras perspectivas; y con qué fiabilidad.
Además, se evaluará si parte de los parámetros se pueden obtener a partir de fuentes públicas cartográficas, sin necesidad de procesar las imágenes.

A continuación se detalla una \textbf{lista de objetivos principales}, que comprenden el \textbf{\gls{MVP}} del trabajo; y una lista de objetivos secundarios de valor añadido.

\subsection{Objetivos principales (MVP)}

Este listado comprende los objetivos planificados para cubrir obligatoriamente, siendo la columna vertebral del proyecto.
\\

\textbf{OP1. Identificar necesidades y requisitos de accesibilidad universal}

Analizar e identificar las necesidades de movilidad y accesibilidad en el espacio público de los entornos urbanos para toda la ciudadanía —teniendo en cuenta la diversidad funcional (física, auditiva, visual, cognitiva, relacional), así como al resto de la población— con el fin de definir los requisitos que garanticen la accesibilidad universal en el entorno urbano, y cumplan el principio de \gls{diseno-para-todos}.
\\

\textbf{OP2. Identificar, analizar y preparar los datos}

Recopilar los conjuntos de datos disponibles de las ciudades de Albacete y Algeciras, analizar los parámetros medidos, seleccionar y categorizar aquellos relevantes según puedan obtenerse de imágenes aéreas y cartografía pública, de imágenes a pie de calle o solo in situ, y preparar los datos para su uso en el entrenamiento y evaluación de los modelos de clasificación.
\\

\textbf{OP3. Identificar y definir las métricas e indicadores de evaluación de la accesibilidad}

Analizar la normativa vigente (estatal, autonómica y local) para identificar el conjunto de indicadores cuantitativos y cualitativos (tales como la anchura libre de paso, o criterios necesarios para plataformas únicas), definir la variable objetivo de la clasificación a partir de las puntuaciones de los datos y establecer las fórmulas matemáticas de las métricas con las que se evaluarán los modelos.
\\

\textbf{OP4. Entrenar modelos de aprendizaje automático de clasificación de accesibilidad}

Modelar y entrenar diferentes modelos de \textit{machine learning}, con diferentes conjuntos de características extraídas de los datos recopilados, diferenciando entre las que se puedan obtener potencialmente de imágenes aéreas o transversales, y las que se obtienen directamente de los datos de campo; comparando los conjuntos aéreo, a pie de calle, la combinación de ambos y el completo, así como un conjunto con variables obtenidas a partir de fuentes cartográficas públicas, con un mismo modelo de referencia.
\\

\textbf{OP5. Validación de los modelos}

Ejecutar el algoritmo sobre los conjuntos de datos de las ciudades piloto, y evaluar de manera cuantitativa la precisión de las predicciones obtenidas por el modelo, cruzando los datos contra los registros previamente calculados sobre el terreno, tanto dentro de cada ciudad como entrenando con una y evaluando en la otra.
\\

\textbf{OP6. Interpretación de los resultados}

Responder a las preguntas de investigación planteadas, interpretando los resultados obtenidos en términos de la accesibilidad de las vías urbanas y la efectividad de la selección de características para imágenes aéreas y cartografía pública, transversales y datos de campo, e identificando qué dimensiones medidas in situ son prescindibles.
\\

\begin{figure}[htbp]
  \centering
  \begin{tikzpicture}[tfm,
      every node/.append style={execute at begin node={\begin{hyphenrules}{nohyphenation}}, execute at end node={\end{hyphenrules}}}, % sin guiones
      obj/.style={arqbase, fill=arq-datos, draw=uoc-azul, text width=40mm, minimum height=19mm,
                  inner sep=4pt, font=\sffamily\scriptsize, align=left, text=tinta},
      objalg/.style={obj, fill=nn-conv, double=nn-conv, double distance=1.2pt},
      num/.style={circle, fill=uoc-azul, text=white, font=\sffamily\scriptsize\bfseries,
                  inner sep=0pt, minimum size=5mm},
    ]
    % ---------- Fondo del MVP
    \fill[uoc-gris] (-2.55, 2.15) rectangle (12.55,-4.65);
    \node[anchor=north west, font=\sffamily\scriptsize\bfseries, text=uoc-azul] at (-2.35, 2.05)
      {Producto mínimo viable (MVP): objetivos principales, en orden de ejecución};

    % ---------- Objetivos
    \node[obj]    (o1) at (0,0)   {\textbf{OP1}\\[1pt] Identificar necesidades y requisitos de accesibilidad universal};
    \node[obj]    (o2) at (5,0)   {\textbf{OP2}\\[1pt] Identificar, analizar y preparar los datos};
    \node[obj]    (o3) at (10,0)  {\textbf{OP3}\\[1pt] Identificar y definir las métricas e indicadores de evaluación de la accesibilidad};
    \node[objalg] (o4) at (0,-3.3) {\textbf{OP4}\\[1pt] Entrenar modelos de aprendizaje automático de clasificación de accesibilidad};
    \node[obj]    (o5) at (5,-3.3) {\textbf{OP5}\\[1pt] Validación de los modelos};
    \node[obj]    (o6) at (10,-3.3){\textbf{OP6}\\[1pt] Interpretación de los resultados};

    % ---------- Números de orden
    \foreach \i in {1,...,6}
      \node[num] at (o\i.north west) {\i};

    % ---------- Flechas en orden
    \draw[flecha] (o1) -- (o2);
    \draw[flecha] (o2) -- (o3);
    \draw[flecha] (o3.south) -- ++(0,-0.62) -| (o4.north);
    \draw[flecha] (o4) -- (o5);
    \draw[flecha] (o5) -- (o6);
  \end{tikzpicture}

  \caption{Objetivos principales del proyecto (MVP).}
  \label{fig:objetivos_mvp}
  \fuente{elaboración propia, con ayuda de Claude.}
\end{figure}


\subsection{Objetivos secundarios}

Representan las metas adicionales o de extensión al trabajo principal del proyecto. Debido a las limitaciones temporales, y al tiempo de dedicación disponible, la consecución de estos objetivos es un valor añadido.
\\

\textbf{OS1. Visualizar las predicciones en un mapa}

Visualizar en un mapa las predicciones de accesibilidad generadas por los modelos entrenados, permitiendo a los usuarios explorar la información obtenida de manera visual e interactiva.
\\

\textbf{OS2. Prueba de concepto con datos remotos reales}

Obtener variables a partir de fuentes públicas (pendientes, anchuras, cruces del callejero) y comprobar si pueden sustituir a la toma de variables in situ.
\\

\subsection{Hipótesis de investigación}
Se plantean las siguientes hipótesis de investigación:

\begin{table}[H]
  \centering
\caption{Hipótesis}
\label{tab:hipotesis}
\small
\renewcommand{\arraystretch}{1.15}
\setlength{\extrarowheight}{1pt}

\begin{tabularx}{\textwidth}{c >{\raggedright\arraybackslash\hsize=1.3\hsize}X >{\raggedright\arraybackslash\hsize=0.7\hsize}X}
\toprule
\cabeceratabla
& \textbf{Hipótesis} & \textbf{Se refuta si…} \\
\midrule
\textbf{H1} & Utilizando únicamente variables que se puedan obtener potencialmente de forma remota (imágenes aéreas y cartografía pública), se clasifica la accesibilidad con una precisión suficiente, similar a cuando se evalúa con todos los parámetros & El rendimiento es claramente inferior al del modelo con todas las variables \\
\midrule
\textbf{H2} & Utilizando únicamente variables que se puedan obtener potencialmente a partir de imágenes a pie de calle, se clasifica la accesibilidad con una precisión suficiente, similar a cuando se evalúa con todos los parámetros & El rendimiento es claramente inferior al del modelo con todas las variables \\
\midrule
\textbf{H3} & Combinar las variables de ambos tipos de imagen (aéreas y cartografía pública, y a pie de calle) mejora el rendimiento de los modelos que utilizan únicamente uno de los dos subconjuntos & No se aprecia mejora suficiente \\
\midrule
\textbf{H4} & El entrenamiento del modelo con los datos de una ciudad es transferible a otra ciudad (es decir, tiene un rendimiento similar) & Hay bastante diferencia de precisión \\
\midrule
\textbf{H5} & Las dimensiones que exigen traslado in situ no son necesarias para obtener un modelo de calidad suficiente & Al excluir las características se detecta un claro descenso en la precisión \\
\bottomrule
\end{tabularx}
\end{table}


\section{Impacto en sostenibilidad, ético-social y de diversidad}
\label{s:etic}

Una vez analizados los Objetivos de Desarrollo Sostenible (ODS) de la Agenda 2030 de la ONU, se enumeran a continuación los impactos positivos, negativos y su mitigación, en relación al presente trabajo fin de máster.

\subsection{Sostenibilidad y huella ecológica}

Los ODS a los que más se contribuye en el proyecto son: \textbf{ODS 11 (\textit{Sustainable Cities and Communities})} y \textbf{ODS 9 (\textit{Industry, Innovation and Infraestructure})}.
\\

En cuanto a los \textbf{impactos positivos}, se establecen:
\begin{itemize}
    \item \textbf{Sustitución de trabajos in situ}: el mecanismo tradicional de auditoría de accesibilidad incluye un desplazamiento continuo de vehículos, para que el personal técnico vaya a cada calle, así como material técnico para poder realizar las mediciones. La utilización de variables extraídas de imágenes reduce de manera significativa los desplazamientos, permitiendo una reducción de las emisiones indirectas CO2.
    \item \textbf{Optimización de los recursos públicos}: además de la reducción de costes de gasolina por los desplazamientos, la posibilidad de ordenar los tramos de calles menos accesibles permite a los ayuntamientos la priorización de obras de reparación en los tramos más críticos, optimizando las intervenciones en los puntos más necesarios.
\end{itemize}

Sin embargo, como \textbf{impacto negativo} se encuentra:
\begin{itemize}
    \item \textbf{Consumo energético por entrenamiento de los modelos de inteligencia artificial}: la utilización de algoritmos de aprendizaje automático va asociada una huella de carbono digital.
\end{itemize}

La \textbf{mitigación} de este punto se basa en el estudio del estado del arte de los modelos de aprendizaje automático para tareas similares, y en caso de utilizar modelos de deep learning, la preferencia del uso de modelos ya preentrenados, utilizando técnicas de \gls{transfer-learning}, lo que hace reducir el tiempo de uso eléctrico y computacional a la hora de entrenar el modelo.

\subsection{Comportamiento ético y responsabilidad social}

El ODS que aplica a este punto es: \textbf{ODS 16 (Peace, Justice and Strong Institutions)}
\\

Como \textbf{impactos positivos}, principalmente en gobernanza abierta y no discriminación, se incluyen:
\begin{itemize}
    \item \textbf{Democratización del dato mediante estándares abiertos}: los resultados de las predicciones del sistema se publican en formatos abiertos (como CSV), por lo que se elimina cualquier dependencia de aplicaciones propietarias de pago, o la necesidad de contratar consultoría tecnológica privada. Al incluir en líneas futuras un portal para que municipios pequeños puedan utilizar la herramienta, se permite la igualdad de condiciones.
    \item \textbf{Transparencia institucional y auditoría social}: la implantación futura propuesta de un catálogo abierto de índice de accesibilidad por municipio permite a la ciudadanía y a los organismos a visibilizar el grado de cumplimiento con la ley de cada localidad, habilitando una auditoría pública.
\end{itemize}

\subsection{Diversidad, género y derechos humanos}

Uno de los objetivos principales del trabajo es el enfoque accesible universal, por lo que los ODS a los que se vincula son \textbf{ODS 10 (Reduced Inequalities)} y \textbf{ODS 5 (Gender Equality)}.
\\

Sobre los \textbf{impactos positivos}, se encuentran:
\begin{itemize}
    \item \textbf{Inclusión y \gls{diversidad-funcional} global}: los beneficios de los planes de accesibilidad no están únicamente pensados para colectivos particulares, sino que están pensados para todo el conjunto diverso de la población. Por lo que no se benefician unicamente en personas con discapacidad permanente, también disfrutan de los planes las personas mayores, personas con discapacidad temporal, familias o el resto de la población.
    \item \textbf{Derecho a la autonomía y disfrute de la ciudad}: el proyecto facilita la implantación de actuaciones de eliminación de barreras en la vía pública y tiene una visión global de la diversidad, por lo que permite una mayor autonomía personal, así como un mejor derecho de disfrute de la ciudad. 
    \item \textbf{Perspectiva de género en la movilidad urbana}: el aumento de autonomía en la diversidad funcional implica una menor necesidad de ayuda para los desplazamientos, algo que impacta positivamente en la habitual mayor carga de cuidados de las mujeres.
\end{itemize}

\section{Enfoque y método seguido}

Utilizando la referencia \textit{Researching Information Systems and Computing} de \textit{Oates}, y analizando la tipología del trabajo, se ha elegido como \textbf{estrategia principal} la de \textbf{investigación basada en diseño}, debido a que se va a realizar la entrega de un sistema informático con modelos de aprendizaje automático para la evaluación de accesibilidad, usando una estrategia adicional de estudio de caso para comparar el sistema con el visor de accesibilidad del ayuntamiento de Albacete.
\\

\subsection{Estrategia principal: investigación basada en diseño (\textit{Design Research})}

Las fases que comprende esta estrategia son:
\begin{enumerate}
    \item \textbf{Identificación del problema (\textit{Awareness})}: identificación del problema grave de las ciudades que no cumplen con criterios de accesibilidad única, análisis de la legislación y buenas prácticas recomendadas por organizaciones accesibles, análisis de los planes de actuación y comprobación de las limitaciones actuales de los procedimientos manuales de verificación de parámetros de accesibilidad, con visitas in situ del personal técnico de los ayuntamientos.
    \item \textbf{Diseño (\textit{Suggestion}):} análisis de literatura científica con el estado del arte del diseño de soluciones para el uso de información tabular de informes de accesibilidad, la reducción de características relevantes, y adicionalmente para el análisis de imágenes aéreas urbanas y extracción de medidas geométricas a partir de datos de cartografía pública. Realización de una propuesta de diseño tecnológica con técnicas de aprendizaje automático, así como un análisis inicial de los conjuntos de datos de Albacete y Algeciras, con diversos modelos para cada conjunto de características.
    \item \textbf{Desarrollo (\textit{Development})}: implementación del sistema informático en Python, incluyendo los modelos y los artefactos web necesarios para cumplir con el diseño planificado.
    \item \textbf{Evaluación (\textit{Evaluation})}: validación cuantitativa de la precisión de los modelos, comparándolo con los datos disponibles de ambas ciudades, lo que se conoce como \gls{ground-truth} en la ciencia de datos. Así como respuesta a las preguntas de investigación planteadas.
    \item \textbf{Conclusión (\textit{Conclusion})}: documentación de los resultados obtenidos, hipótesis confirmadas o refutadas, listado de variables recomendadas para futuros estudios de extracción de imágenes; aprendizajes y dificultades encontradas durante la realización del trabajo, así como enumeración de futuras líneas de ampliación.
\end{enumerate}

\begin{figure}[htbp]
  \centering
  \begin{tikzpicture}[tfm,
      every node/.append style={execute at begin node={\begin{hyphenrules}{nohyphenation}}, execute at end node={\end{hyphenrules}}}, % sin guiones
      fase/.style={arqbase, fill=arq-datos, draw=uoc-azul, text width=46mm, minimum height=22mm,
                   inner sep=4pt, font=\sffamily\scriptsize, align=left, text=tinta},
      fasealg/.style={fase, fill=nn-conv, double=nn-conv, double distance=1.2pt},
      num/.style={circle, fill=uoc-azul, text=white, font=\sffamily\scriptsize\bfseries,
                  inner sep=0pt, minimum size=5mm},
    ]
    % ---------- Etapas dispuestas en círculo (sentido horario)
    \node[fase]    (f1) at (0, 3.7)   {\textbf{Identificación del problema}\\[1pt]
      Descubrir el objetivo del proyecto a partir de la lentitud de los planes de accesibilidad, que exigen inspecciones \emph{in situ}.};
    \node[fase]    (f2) at (5.3, 1.2) {\textbf{Diseño}\\[1pt]
      Análisis de las leyes y del estado del arte para elegir las características de medida de la accesibilidad, y análisis de los datos de Albacete y Algeciras.};
    \node[fase] (f3) at (3.3,-3.1) {\textbf{Desarrollo}\\[1pt]
      Implementación de los modelos y de los sistemas.};
    \node[fase]    (f4) at (-3.3,-3.1) {\textbf{Evaluación}\\[1pt]
      Validación de la precisión de los modelos y comparativa entre modelos y entre ciudades.};
    \node[fase]    (f5) at (-5.3, 1.2) {\textbf{Conclusión}\\[1pt]
      Respuesta a las preguntas de investigación, documentación de los resultados y líneas futuras.};

    % ---------- Números de etapa
    \foreach \i in {1,...,5}
      \node[num] at (f\i.north west) {\i};

    % ---------- Flechas del ciclo
    \draw[flecha] (f1.east)  to[bend left=22] (f2.north);
    \draw[flecha] (f2.south) to[bend left=18] (f3.north);
    \draw[flecha] (f3.west)  -- (f4.east);
    \draw[flecha] (f4.north) to[bend left=18] (f5.south);
    \draw[flecha, dashed] (f5.north) to[bend left=22] node[pos=0.5, above left, etiqueta] {nuevo ciclo} (f1.west);

    % ---------- Centro
    \node[align=center, font=\sffamily\scriptsize, text=uoc-azul] at (0,0.1)
      {\textbf{Investigación}\\\textbf{basada en diseño}\\[1pt]\textcolor{tinta-suave}{(\emph{Design Research})}};
  \end{tikzpicture}

  \caption{Ciclo de vida de la metodología de desarrollo del trabajo.}
  \label{fig:metodologia_ciclo}
  \fuente{elaboración propia, con ayuda de Claude.}
\end{figure}


\subsection{Estrategia adicional: caso de estudio (\textit{Case Study})}

Tomando como base el conjunto de datos de Albacete y Algeciras, se puede realizar una evaluación del comportamiento del sistema.

De esta manera, se podrá demostrar la calidad del sistema, y servirá para demostrar el uso real de la herramienta en los ayuntamientos.

\subsection{Metodología de ingeniería de software}
Se plantea la ejecución del diseño e implementación de código como un proyecto de ingeniería de software, con tecnologías de codigo abierto, herramientas de inteligencia artificial, y de gestión de proyectos.

\begin{enumerate}
    \item Metodología
        \begin{itemize}
            \item \textbf{Desarrollo ágil con metodología \gls{kanban}:} organización del trabajo utilizando tableros \gls{kanban} en \gls{github} projects.
            \item \textbf{Principios de código limpio y control de versiones:} implementación del código siguiendo los principios de \textit{Clean Code}, utilizando control de versiones con el flujo de \textit{Git Flow}.
        \end{itemize}
    \item Stack tecnológico
        \begin{itemize}
            \item \textbf{Lenguaje de programación:} Python, el lenguaje estándar para los trabajos de ciencia de datos.
            \item \textbf{Librerías de \textit{deep learning} y visión por computador:} uso de librerías de Python como PyTorch o Torchvision.
            \item \textbf{Librerías de tratamiento espacial de datos:} utilización de librerías como GeoPandas, NumPy, exportación en formatos abiertos (GeoJson/CSV), y tecnologías estándar de desarrollo web (HTML5, CSS, Javascript).
        \end{itemize}
    \item Recursos hardware
        \begin{itemize}
            \item \textbf{Entorno con \gls{GPU}:} se necesita utilizar un entorno con aceleración gráfica por hardware (\gls{GPU}), ya sea en entorno local o remoto, para poder entrenar los modelos.
        \end{itemize}
    \item Herramientas
        \begin{itemize}
            \item \textbf{Procesador de textos:} se utilizará Overleaf como plataforma de procesador de textos LaTeX para la redacción de la memoria.
            \item \textbf{Repositorio de código:} GitHub como repositorio de código fuente y control de versiones.
            \item \textbf{Seguimiento ágil del proyecto:} se utilizará GitHub Projects para controlar el seguimiento del trabajo, utilizando épicas e historias de usuario.
        \end{itemize}
    \item Plataformas de formación
        \begin{itemize}
            \item Se utilizarán \textbf{plataformas de aprendizaje online} como DataCamp o Coursera para consolidar y refrescar los conceptos de Python y sus librerías de ciencia de datos con impacto en el proyecto.
        \end{itemize}    
    \item Utilización de la inteligencia artificial  
\end{enumerate}


\section{Planificación del trabajo}

Se ha realizado una \textbf{planificación en base a la carga de 12 créditos ECTS}, equivalentes a 300 horas de dedicación, distribuyendo la carga de trabajo en las 15 semanas lectivas disponibles.

La estructuración del trabajo se ha realizado utilizando \textbf{paquetes de trabajo (WP)}, que se han basado principalmente en las entregas establecidas en las pruebas de integración continua de la UOC, incluyendo una serie de \textbf{hitos} para reflejar la evolución incremental del trabajo tanto en el estudio del estado del arte y en el diseño e implementación del mismo.

Además, se han incluido paquetes de trabajo para los \textbf{objetivos secundarios}, planificados inicialmente como opcionales.

\subsection{Calendario del trabajo}

En primer lugar se presenta el \textbf{cronograma del proyecto}, incluyendo tanto los paquetes de trabajo como hilos de entrega, se detalla en el siguiente diagrama de Gantt. En color azul se encuentran las entregas principales, mientras que en color verde se indican las entregas secundarias.

\begin{figure}[H]
  \centering
  \begin{ganttchart}[
      tfmgantt,                          % estilo del sistema (tfm-estilo.sty)
      time slot format=isodate,
      x unit=1.0mm,                      % 120 días en el ancho de página
      y unit chart=0.42cm, y unit title=0.5cm,
      vgrid={*6{draw=none}, *1{draw=rejilla, line width=\trazofino}},  % una línea por semana
      hgrid={*1{draw=rejilla, line width=\trazofino}},
      canvas/.append style={draw=filete, line width=\trazofino},
      bar height=0.55, milestone height=0.6, group peaks height=0.2,
      milestone left shift=-1.5, milestone right shift=1.5,  % rombo de ~4 mm
      bar label font=\sffamily\scriptsize,
      group label font=\sffamily\scriptsize\bfseries\color{uoc-azul},
      milestone label font=\sffamily\scriptsize\itshape\color{tinta-suave},
      title label font=\sffamily\scriptsize\bfseries\color{uoc-azul},
      secundaria/.style={bar/.append style={fill=papel, draw=uoc-azul, dashed}},
      descanso/.style={bar/.append style={fill=uoc-gris, draw=tinta-suave, line width=\trazofino,
                        pattern=north east lines, pattern color=tinta-suave},
                       bar label font=\sffamily\scriptsize\itshape\color{tinta-suave}},
    ]{2026-09-24}{2027-01-21}

    % Escala de tiempo (en castellano; los números son días de cada mes dentro del rango)
    \gantttitle{2026}{99} \gantttitle{2027}{21} \\
    \gantttitle{sep.}{7} \gantttitle{oct.}{31} \gantttitle{nov.}{30}
    \gantttitle{dic.}{31} \gantttitle{ene.}{21} \\

    % WP1
    \ganttbar{WP1 Definición y objetivos}{2026-09-24}{2026-10-04} \\
    \ganttmilestone{Entrega PEC1}{2026-10-04} \\

    % WP2
    \ganttgroup{WP2 Contexto}{2026-10-05}{2026-10-25} \\
    \ganttbar{WP2.1 Análisis legal}{2026-10-05}{2026-10-12} \\
    \ganttbar{WP2.2 Estado del arte}{2026-10-13}{2026-10-25} \\
    \ganttmilestone{Entrega PEC2}{2026-10-25} \\

    % WP3
    \ganttgroup{WP3 Desarrollo}{2026-10-19}{2026-12-15} \\
    \ganttbar{WP3.1 Preparación de datos}{2026-10-19}{2026-11-03} \\
    \ganttbar{WP3.2 Métricas y características}{2026-10-26}{2026-11-10} \\
    \ganttbar{WP3.3 Entrenamiento de modelos}{2026-11-03}{2026-11-24} \\
    \ganttbar{WP3.4 Comparativa y ajuste}{2026-11-11}{2026-11-24} \\
    \ganttbar{WP3.5 Validación e iteración}{2026-11-13}{2026-12-05} \\
    \ganttbar{WP3.6 Interpretación de resultados}{2026-11-25}{2026-12-05} \\
    \ganttbar[descanso]{Descanso}{2026-12-02}{2026-12-07} \\
    \ganttbar[secundaria]{WP3.7 (OS1) Integración con el callejero}{2026-12-08}{2026-12-11} \\
    \ganttbar[secundaria]{WP3.8 (OS1) Visor GIS e interfaz}{2026-12-10}{2026-12-14} \\
    \ganttbar[secundaria]{WP3.9 (OS2) POC con imágenes}{2026-12-08}{2026-12-15} \\
    \ganttmilestone{Entrega PEC3}{2026-12-15} \\

    % WP4
    \ganttbar{WP4 Redacción memoria}{2026-12-16}{2026-12-27} \\
    \ganttmilestone{Entrega PEC4}{2026-12-27} \\

    % WP5 y WP6
    \ganttbar{WP5 Presentación vídeo}{2026-12-28}{2027-01-05} \\
    \ganttmilestone{Entrega PEC5}{2027-01-05} \\
    \ganttbar{WP6 Defensa oral}{2027-01-06}{2027-01-21} \\
    \ganttmilestone{Entrega PEC6 / Final}{2027-01-21}
  \end{ganttchart}

  \vspace{3mm}
  % Leyenda
  \begin{tikzpicture}[font=\sffamily\scriptsize, text=tinta]
    \draw[uoc-azul, fill=uoc-azul] (0,-0.06) rectangle (0.5,0.06);
    \node[anchor=west] at (0.6,0) {Fase};
    \draw[uoc-azul, fill=nn-conv, line width=\trazofino] (1.6,-0.1) rectangle (2.1,0.1);
    \node[anchor=west] at (2.2,0) {Tarea};
    \draw[uoc-azul, fill=papel, dashed, line width=\trazofino] (3.3,-0.1) rectangle (3.8,0.1);
    \node[anchor=west] at (3.9,0) {Tarea secundaria};
    \draw[tinta-suave, fill=uoc-gris, line width=\trazofino, pattern=north east lines, pattern color=tinta-suave] (6.3,-0.1) rectangle (6.8,0.1);
    \node[anchor=west] at (6.9,0) {Descanso};
    \draw[acc-no-accesible, fill=acc-deficiente] (8.35,0) -- ++(0.1,0.1) -- ++(0.1,-0.1) -- ++(-0.1,-0.1) -- cycle;
    \node[anchor=west] at (8.6,0) {Hito (entrega de PEC)};
  \end{tikzpicture}

  \caption{Diagrama de Gantt del proyecto.}
  \label{fig:gantt_tfm}
\end{figure}


En segundo lugar, se muestra en la siguiente tabla la división de las tareas a realizar durante el trabajo, descompuestas en \textbf{6 paquetes de trabajo principales}, indicando el número de créditos ETCS y horas, y a qué entrega de evaluación continua se corresponden. Para el hito de implementación (WP3), se ha realizado una \textbf{división inicial en subfases}, que podría verse actualizada una vez se consolide la investigación del estado del arte (WP2).

\begin{table}[H]
  \centering
\caption{Planificación del proyecto y distribución de carga ECTS}
\label{tab:planificacion_tfm}
\scriptsize
\setlength{\extrarowheight}{1pt}

\begin{tabularx}{\textwidth}{l X c c c}
\toprule
\cabeceratabla
\textbf{WP} & \textbf{Tarea y descripción detallada} & \textbf{Carga} & \textbf{Fechas} & \textbf{UOC} \\
\midrule
    WP1 & \textbf{Definición y objetivos:} planteamiento inicial, problemática, alineación de objetivos, metodología y planificación temporal. & \makecell{0.6 ECTS \\ (15 h)} & 24 sep - 04 oct & PEC1 \\
\midrule
    WP2.1 & \textbf{Análisis legal:} estudio de legislación y normativa sobre accesibilidad e infraestructura urbana, extrayendo requisitos. & \makecell{0.6 ECTS \\ (15 h)} & 05 oct - 12 oct & PEC2 \\
    WP2.2 & \textbf{Estado del arte:} revisión de literatura sobre evaluación de la accesibilidad urbana, aprendizaje automático con datos tabulares, validación con datos espaciales y visión por computador. & \makecell{1.0 ECTS \\ (25 h)} & 13 oct - 25 oct & PEC2 \\
\midrule
    WP3.1 & \textbf{Preparación de datos:} recopilación de los datos de Albacete y Algeciras, análisis de los parámetros medidos, selección y categorización de los relevantes, limpieza y preparación para el entrenamiento y la evaluación de los modelos. & \makecell{0.8 ECTS \\ (20 h)} & 19 oct - 03 nov & PEC3 \\
    WP3.2 & \textbf{Métricas y características:} definición de los indicadores y fórmulas de evaluación a partir de la normativa, y construcción de los conjuntos de características (aéreas, transversales y de campo). & \makecell{0.6 ECTS \\ (15 h)} & 26 oct - 10 nov & PEC3 \\
    WP3.3 & \textbf{Entrenamiento de modelos:} entrenamiento de modelos de aprendizaje automático de clasificación de accesibilidad con los distintos conjuntos de características. & \makecell{1.2 ECTS \\ (30 h)} & 03 nov - 24 nov & PEC3 \\
    WP3.4 & \textbf{Comparativa y ajuste:} comparación de los modelos entre sí y entre los distintos conjuntos de características, y ajuste de hiperparámetros. & \makecell{0.8 ECTS \\ (20 h)} & 11 nov - 24 nov & PEC3 \\
    WP3.5 & \textbf{Validación e iteración:} evaluación frente al conjunto de prueba, análisis de errores y ajuste continuo de hiperparámetros. & \makecell{0.6 ECTS \\ (15 h)} & 13 nov - 05 dic & PEC3 \\
    WP3.6 & \textbf{Interpretación de resultados:} análisis e interpretación de los resultados para responder a las preguntas de investigación, en términos de la accesibilidad de las vías urbanas y de la efectividad de la selección de características. & \makecell{0.6 ECTS \\ (15 h)} & 25 nov - 05 dic & PEC3 \\
\rowcolor{gray!15}
    --- & \textbf{Descanso:} pausa planificada sin carga lectiva ni entregas asignadas. & \makecell{0.0 ECTS \\ (0 h)} & 02 dic - 07 dic & --- \\
\rowcolor{gray!6}
    \multicolumn{5}{l}{\textit{Objetivos secundarios (de valor añadido, condicionados al tiempo disponible y a la evolución del proyecto)}} \\
    WP3.7 (OS1) & \textbf{Integración con el callejero:} realización de la integración del modelo con los datos geográficos del callejero de una ciudad, para poder incorporar los resultados del análisis al visor. & \makecell{0.4 ECTS \\ (10 h)} & 08 dic - 11 dic & PEC3 \\
    WP3.8 (OS1) & \textbf{Visor GIS e interfaz:} desarrollo de un visor geográfico interactivo para la representación gráfica sobre mapas de la red urbana. & \makecell{0.4 ECTS \\ (10 h)} & 10 dic - 14 dic & PEC3 \\
    WP3.9 (OS2) & \textbf{Prueba de concepto con datos remotos reales:} extracción de información de fuentes cartográficas públicas y evaluación de la posibilidad de obtener los datos necesarios para las predicciones del modelo de clasificación. & \makecell{0.8 ECTS \\ (20 h)} & 08 dic - 15 dic & PEC3 \\
\midrule
    WP4 & \textbf{Redacción memoria:} consolidación y redacción de la memoria final del TFM, estructurando metodología y resultados. & \makecell{1.6 ECTS \\ (40 h)} & 16 dic - 27 dic & PEC4 \\
\midrule
    WP5 & \textbf{Presentación vídeo:} elaboración del guion, soporte visual y grabación del recurso audiovisual requerido para la defensa. & \makecell{1.0 ECTS \\ (25 h)} & 28 dic - 05 ene & PEC5 \\
\midrule
    WP6 & \textbf{Defensa oral:} ensayo, preparación de respuestas a posibles objeciones del tribunal y exposición oral final del proyecto. & \makecell{1.0 ECTS \\ (25 h)} & 06 ene - 21 ene & PEC6 \\
\bottomrule
\end{tabularx}
\end{table}


\subsection{Identificación de riesgos}
Con el fin de asegurar el cumplimiento de la planificación, así como la consecución de los hitos de entrega académicos, se ha realizado una matriz de evaluación de posibles riesgos que pueden tener impacto en la elaboración del trabajo, incluyendo qué probabilidad tienen de producirse, qué impacto pueden tener en el proyecto, y el plan de contingencia a seguir si se producen.

Los principales riesgos identificados han estado centrados en la incentidumbre de la calidad de entrenamiento y predicción de los modelos de aprendizaje profundo, así como la priorización de objetivos principales frente a los secundarios.

\begin{table}[H]
  \centering
  \caption{Matriz de identificación de riesgos y planes de contingencia}
  \label{tab:identificacion_riesgos}
  \scriptsize
  \setlength{\extrarowheight}{3pt}
  
  \begin{tabularx}{\textwidth}{>{\raggedright\arraybackslash}p{2.8cm} c >{\raggedright\arraybackslash}X >{\raggedright\arraybackslash}X}
    \toprule
    \cabeceratabla
    \textbf{Factor de riesgo} & \textbf{Prob. / Impacto} & \textbf{Impacto en el proyecto} & \textbf{Plan de contingencia} \\
    \midrule
    
    \textbf{Dificultad en el análisis de los conjuntos de datos de Albacete o Algeciras} & 
    Media/Alto & 
    Constituye uno de los principales riesgos del proyecto al ser el punto de inicio del resto del sistema. Problemas posibles: las variables no tienen descripción clara por lo que puede dificultar su interpretación y análisis, así como la dificultad de comparar ambos conjuntos de datos. & 
    Análisis inicial de los datos en la segunda entrega del proyecto, y consulta con el tutor o expertos que han generado la información de los conjuntos de datos.
    Generación de variables derivadas para facilitar el análisis y comparación de los conjuntos de datos. \\
    \midrule
    
    \textbf{Tiempos elevados de entrenamiento de modelos} & 
    Media/Medio & 
    Los trabajos de ciencia de datos se basan en la experimentación y prueba continua, por lo que el tiempo de entrenamiento constituye una variable importante en la consecución de objetivos. Con tiempos de entrenamiento elevados, se reduce el número de iteraciones que se pueden realizar y validar. & 
    Priorización del uso de modelos con mejor optimización y en caso de utilizar deep learning, usar \textit{transfer learning} sobre redes preentrenadas basándose en el estado del arte analizado y empleo de aceleración por GPU. \\
    \midrule
    
    \textbf{No llegar a tiempo para completar los objetivos secundarios} & 
    Media/Bajo & 
    En caso de que la planificación de los objetivos principales no se cumpla a tiempo, hay un riesgo de que se extienda el tiempo dedicado a estos objetivos, influyendo en el tiempo disponible para realizar los objetivos secundarios del proyecto. & 
    Dar prioridad a los hitos principales, extendiendo la dedicación en el cronograma y reevaluando cada semana el estado de los hitos secundarios en la planificación. \\
    \midrule

    \textbf{Dificultad en obtener datos de cartografía pública} &
    Alta/Medio &
    La obtención de datos de cartografía pública que sean útiles para los modelos, así como la precisión de los datos es un riesgo que puede darse, así como enlazar los datos extraídos con las variables disponibles. &
    No añadir todas las fuentes cartográficas a la vez, y realizar previamente un análisis de viabilidad. \\
    
    \textbf{Indisponibilidad personal por vacaciones planificadas} & 
    Alta/Medio & 
    Se ha incluido en el calendario un período de descanso de una semana (del 2 al 7 de diciembre) por un período vacacional planificado, lo que puede afectar al tiempo disponible para la realización del proyecto. & 
    Adelanto de la fase implementación (WP3) respecto a lo indicado por el calendario de entregas de la UOC. \\
    
    \bottomrule
  \end{tabularx}
\end{table}


\section{Breve sumario de productos obtenidos}

Los productos que se obtendrán tras la ejecución del proyecto se detallan a continuación, distinguiendo los entregablos del MVP de los secundarios:

\begin{table}[H]
  \centering
\caption{Productos del proyecto.}
\label{tab:productos2}
\small
\renewcommand{\arraystretch}{1.15}
\begin{tabularx}{\textwidth}{l >{\raggedright\arraybackslash}X l l}
\toprule
\cabeceratabla
    Código & Producto & Formato & WP \\
\midrule
\multicolumn{4}{l}{\sffamily\footnotesize\bfseries\color{uoc-azul}Producto mínimo viable (MVP)} \\
    P1 & Conjunto de datos limpio y preparado de las vías urbanas de Albacete y Algeciras, con su diccionario de datos. & Dataset & WP3.1 \\
    P2 & Repositorio en GitHub con el código de preparación de datos, selección de características y entrenamiento de los modelos de clasificación de accesibilidad. & Código & WP3.2--3.3 \\
    P3 & Comparativa de modelos y de conjuntos de variables (aéreas, a pie de calle y de campo), con evaluación de los mismos y rendimiento entre ciudades. & Resultados y código & WP3.4--3.5 \\
    P4 & Análisis e interpretación de los resultados frente a las hipótesis, indicando qué variables sirven para poder clasificar la accesibilidad. & Informe & WP3.6 \\
\midrule
\multicolumn{4}{l}{\sffamily\footnotesize\bfseries\color{uoc-azul}Productos secundarios\textsuperscript{*}} \\
    P5 & Capa de calles clasificadas integrada en el callejero municipal. & Capa GIS & WP3.7 \\
    P6 & Visor GIS interactivo con los resultados. & Aplicación web & WP3.8 \\
    P7 & Prueba de concepto de extracción de variables a partir de datos cartográficos. & Código e informe & WP3.9 \\
\midrule
\multicolumn{4}{l}{\sffamily\footnotesize\bfseries\color{uoc-azul}Entregables académicos} \\
    P8  & Memoria del TFM. & PDF & WP4 \\
    P9  & Presentación en vídeo. & Vídeo & WP5 \\
\bottomrule
\end{tabularx}
\notatabla{\textsuperscript{*}\,Productos secundarios: se desarrollarán si el calendario lo permite.}
\end{table}

\section{Breve descripción de los otros capítulos de la memoria}

La estructura de la memoria tendrá los siguientes capítulos, coincidiendo en orden con las entregas estipuladas en el portal de la universidad:

\begin{enumerate}
    \item \textbf{Introducción y contexto del proyecto:} define el contexto y motivación del proyecto, una introducción a la normativa legal, el listado de objetivos, la metodología y la planificación del trabajo.
    \item \textbf{Características de vías accesibles. Marco legal y normativo:} define la legislación vigente en relación a las ciudades accesibles universales, así como las características de las vías accesibles en base a la ley, que servirán de base para los algoritmos.
    \item \textbf{Estado del arte:} contiene el análisis del estado el arte científico en cuanto a modelos de aprendizaje automático y algoritmos utilizados en problemas similares.
    \item \textbf{Arquitectura del sistema:} describe la arquitectura tecnológica del sistema, incluyendo detalles de cada uno de los módulos del mismo: algoritmos de \textit{deep learning}, de visión computacional, y clasificación de accesibilidad de calles.
    \item \textbf{Diseño e implementación de los algoritmos:} contiene las decisiones de diseño de los sistemas, las librerías y tecnológícas utilizadas, y los detalles y decisiones de implementación.
    \item \textbf{Resultados y evaluación:} detalla los resultados de las predicciones realizadas, lista las rutas a los sistemas y conjuntos de datos, y analiza la evaluación de la calidad del trabajo.
    \item \textbf{Conclusiones y trabajos futuros:} contiene las conclusiones finales del estudiante y el listado de líneas futuras de trabajos.
\end{enumerate}

\section{Declaración sobre el uso de inteligencia artificial generativa}

En primer lugar, se listan los casos de uso de la inteligencia artificial generativa durante el desarrollo del trabajo:

\begin{itemize}
    \item \textbf{Facilitación de ingesta de legislación y literatura científica:} se ha usado NotebookLM para poder analizar y comprender mejor textos legales y científicos y generando artefactos como mapas mentales o infografías conceptuales, o para obtener las secciones más interesantes a leer de las fuentes seleccionadas.
    \item \textbf{Recursos gráficos para la memoria:} con herramientas de inteligencia artificial generativa para incluir imágenes o esquemas que faciliten la lectura de la memoria y la enriquezcan.
    \begin{itemize}
        \item En primer lugar, se ha generado un sistema de diseño utilizando Claude Design, a partir de la plantilla de la memoria de la UOC, que incluye una paleta de colores, tipografías y componentes (gráficos, tablas, diagramas...) que permiten enriquecer visualmente la memoria y mantener un estilo coherente en todo el documento.
    \end{itemize}
    \item \textbf{Asistencia en formateo de textos en LaTeX:} como apoyo técnico para poder transformar textos escritos en textos formateados para LaTeX, o maquetación de tablas o diagramas.
    \item \textbf{Asistencia programación y refactorización de código:} se utilizarán asistentes de código como Claude o Copilot como ayuda en la escritura de \textit{scripts} en Python.
    \item \textbf{Diseño UI/UX y estilos del portal web:} con el fin de ofrecer una página web más atractiva para el usuario, se emplearán herramientas de inteligencia artificial que den soporte con el diseño y estilos de la web.
\end{itemize}

Para ver detalles de los prompt específicos usados, se puede consultar el \nameref{anx:ia}.